% Generated draft: requires mathematical and source review.
\documentclass[11pt]{article}
\usepackage[margin=25mm]{geometry}
\usepackage{fontspec}
\setmainfont{Latin Modern Roman}
\usepackage{amsmath,amssymb,mathtools,unicode-math}
\setmathfont{Latin Modern Math}
\usepackage{enumitem,needspace,xurl,hyperref}
\setlength{\parindent}{0pt}
\setlength{\parskip}{5pt}
\setlength{\emergencystretch}{3em}
\newcommand{\N}{\mathbb{N}}
\newcommand{\Z}{\mathbb{Z}}
\newcommand{\Q}{\mathbb{Q}}
\newcommand{\R}{\mathbb{R}}
\newcommand{\C}{\mathbb{C}}
\newcommand{\F}{\mathbb{F}}
\newcommand{\A}{\mathbb{A}}
\newcommand{\PP}{\mathbb{P}}
\newcommand{\E}{\mathbb{E}}
\newcommand{\eps}{\varepsilon}
\newcommand{\dd}{\,\mathrm{d}}
\begin{document}
\section*{UnsolvedMath ID: 30005034}
\textbf{Source problem code:} catalogue entry
\subsection{Definitions and mathematical statement}
A homological filling function measures the size of chains needed to fill cycles of bounded size. More precisely, after fixing a coefficient ring, a degree \(d\), a chain complex \((C_*,\partial_*)\), and chain norms, the corresponding filling expression is
\[
\operatorname{Fill}(z)=\inf\{\|b\|:b\in C_{d+1},\ \partial_{d+1}b=z\},\qquad
F(n)=\sup\{\operatorname{Fill}(z):z\in\ker\partial_d,\ \|z\|\le n\}.
\]
Here an unfillable cycle has filling value \(+\infty\). Defining a group invariant by this expression requires suitable hypotheses on the complex and its relationship to the group. The supplied record does not specify the degree, coefficients, norms, group models, or finiteness hypotheses, so these data remain unresolved.

For functions \(f,g:\mathbb N\to[0,+\infty]\), a customary coarse comparison is
\[
f\preceq g\quad\Longleftrightarrow\quad
\exists C\in\mathbb Z_{\ge1}\ \forall n\in\mathbb N,\quad
f(n)\le Cg(Cn+C)+Cn+C.
\]
Write \(f\simeq g\) when both \(f\preceq g\) and \(g\preceq f\). This is an explicit candidate meaning of equivalence; the supplied passage does not define its intended comparison relation.

\textbf{Question, conditional on resolving these conventions.} Let \(\mathcal G\) denote the intended class of groups with homological filling functions \(F_G\). Is the collection
\[
\{[F_G]_{\simeq}:G\in\mathcal G\}
\]
uncountable? As a proposed first step, for arbitrary admissible proper subsets \(S,T\) of the intended ambient set, is
\[
F_{\operatorname{GL}(S)}\simeq F_{\operatorname{GL}(T)}?
\]
The notation \(\operatorname{GL}(S)\) is retained exactly as supplied. Neither this group construction, its possible additional parameters, the ambient set for \(S,T\), nor their admissibility conditions are defined in the record. The qualification that the subsets are proper comes from the supplied literature assessment. Consequently, the displayed questions describe the recorded targets but do not constitute a fully specified formulation until these omissions are resolved. No logical implication between the two questions is asserted.
\subsection{Short English statement}
The problem asks whether groups exhibit uncountably many distinct homological filling growth types, and proposes comparing the functions of the groups denoted \(\operatorname{GL}(S)\) and \(\operatorname{GL}(T)\) as a first step. The group construction and essential filling conventions are missing, so the formulation remains conditional and requires human mathematical review.
\subsection{Sources}
\begin{itemize}
\item Supplied local UnsolvedMath problem registry, permanent ID \(30005034\), title ``Equivalence Classes of Homological Filling Functions,'' source problem code \texttt{catalogue entry}. No individual proposer is identified. The supplied proposed year is 2022.
\item Supplied authoritative report citation: \emph{Geometric Structures in Group Theory (2023)}, Oberwolfach Reports, pp.\ 517--576, DOI \texttt{10.4171/owr/2022/11}. \url{https://doi.org/10.4171/owr/2022/11}. The report title year and DOI are reproduced as supplied.
\item Supplied literature source: \emph{Homological Dehn functions of groups of type FP2}. \url{https://arxiv.org/abs/2012.00730}.
\item Supplied literature source: ``Oberwolfach report.'' \url{https://doi.org/10.4171/owr/2022/11}.
\item Literature assessment dated exactly \texttt{2026-08-22}: the supplied curation labels the problem open and reports that the cited work establishes broad upper and lower comparisons and equality in special examples, while equivalence for arbitrary proper subsets \(S,T\) remains an open question. This is supplied dated curation, not an independently verified literature assessment.
\item The registry labels the statement \texttt{corrected\_verified} and reports that the official report was checked directly and definitions restored. These are supplied provenance claims; the present record does not contain the definitions needed to substantiate a complete formulation, and no source was independently visited for this draft.
\end{itemize}
\end{document}
